\subsection{Fiber-Induced Structure and the Open Gravitational Selection Problem}
  \label{sec:necessity-noninjective}

  This section separates a proved restriction on fiber-induced structure from open gravitational constructions.
  Injectivity removes vertical re-identification symmetries and, in the principal-bundle setting, fiber holonomy.
  It does not exclude an effective metric or diffeomorphism covariance on the base.
  A bounded relaxation flux is retained as a candidate selection principle, not a proved necessity for gravity.
  Appendix~\ref{subsec:injective-holonomy} supplies the fiber result;
  Appendix~\ref{subsec:IR} records the missing operator-level construction for the infrared action.

  \subsubsection*{Operational definitions and standing assumptions}
    \label{subsec:necessity-defs}

    We work with a microscopic configuration space $\mathcal{C}_\chi$ endowed
    with an admissibility rule $F$ encoding the structural constraints of~$\chi$
    (Section~\ref{sec:definition-and-fundamental-properties-of-the-chi-field}).
    The space of effective observables is $\mathcal{M}$.
    A projection (coarse-graining) map $\Pi : \mathcal{C}_\chi \to \mathcal{M}$ defines the
    effective state $y = \Pi(x)$ for each substrate configuration $x \in \mathcal{C}_\chi$.

    \paragraph{Emergence.}
      An effective quantity $Q_{\mathrm{eff}}$ is called \emph{emergent} if and only if
    \begin{equation}
      Q_{\mathrm{eff}} = Q(\Pi, F, x),
      \end{equation}
      determined solely by the admissibility rule $F$ and the projection $\Pi$.
      No independent free parameter carrying $Q$ is introduced at the effective level.

    \paragraph{Ontological uniqueness.}
      The complete substrate is $(\mathcal{C}_\chi, F)$.
      All effective structures, metric, connection, mass, and charge, are derived from the pair $(F, \Pi)$ only.
      No fundamental geometric field is postulated independently of this pair.

    \paragraph{Observed gravity.}
      We assume the effective description admits a Lorentzian geometric formulation satisfying:
    \begin{enumerate}[label=(D\arabic*)]
  \item \label{D1}
  Effective locality.
  Field equations are local on $\mathcal{M}$.
  \item \label{D2}
  Covariance.
  Diffeomorphism invariance as an invariance of description, not of the substrate.
  \item \label{D3}
  Infrared dominance of a low-derivative action.
  \item \label{D4}
  Universal coupling to energy-momentum.
      \end{enumerate}

    \paragraph{Hypothesis H (projectability).}
      The effective geometry depends on a substrate configuration only through its projected state:
    \begin{equation}
      g_{\mathrm{eff}}(x) = g(\Pi(x)),
      \qquad y=\Pi(x).
      \label{eq:hypothesis-H}
    \end{equation}
      Thus configurations in the same fiber have the same effective geometry.
      This factorization does not prohibit writing a pull-back quantity as a function of $x$.
      If $\Pi$ is bijective onto its image, any quantity on that image can be written using $\Pi^{-1}$;
      Hypothesis~H alone then supplies no obstruction to a derived metric.

  \subsubsection*{Injectivity removes vertical structure}
    \label{subsec:lemma1}

    \begin{lemma}[Injectivity removes fiber re-identification]
      \label{lem:injectivity-no-gravity}
      Let $\Pi : \mathcal{C}_\chi \to \mathcal{M}$ be injective.
      Every transformation $\phi$ satisfying $\Pi\circ\phi=\Pi$ is the identity.
      Every nonempty fiber is a singleton, so its counting multiplicity has zero logarithmic gradient
      wherever that gradient is defined.
      For an admissible principal bundle in the sense of Definition~\ref{def:admissible-projection},
      the structure group and holonomy are trivial and the connection curvature vanishes.
    \end{lemma}

    \begin{proof}
      From $\Pi(\phi(x))=\Pi(x)$ and injectivity, $\phi(x)=x$ for every $x$.
      On the image of $\Pi$, $|\Pi^{-1}(y)|=1$ and $\log|\Pi^{-1}(y)|=0$.
      In the principal-bundle case the structure group acts freely and transitively on each singleton fiber,
      so it is the trivial group.
      Its Lie algebra is zero and its holonomy is a subgroup of the trivial group.
      These are precisely the conclusions of Theorem~\ref{thm:injective-trivial-holonomy}
      and Corollary~\ref{cor:no-gauge}.
    \end{proof}

    This result concerns structure induced by the fibers of $\Pi$, not all structures derivable from $(F,\Pi)$.
    In particular, the Hessian expression in Eq.~\eqref{eq:eff-metric-hessian} uses no fiber connection.
    Injectivity supplies no reason for such a projectable metric, if defined and nondegenerate, to be flat.
    The existence and dynamics of that metric require their own construction.

  \subsubsection*{Descriptive quotients and their domains}
    \label{subsec:lemma2}

    \begin{remark}[Descriptive redundancy need not be a substrate fiber]
      \label{lem:redundancy-quotient}
      A quotient identifying distinct representatives is a non-injective map on its domain of representatives.
      In general relativity, for example, diffeomorphisms identify metric representatives on a fixed manifold.
      That quotient is on the space of effective descriptions, not automatically on $\mathcal{C}_\chi$.
      Identifying it with the substrate projection $\Pi$ would require an additional construction.
      Moreover, a vertical symmetry satisfying $\Pi\circ\phi=\Pi$ fixes every base state.
      Such symmetries cannot implement a diffeomorphism that moves base points.
      Fiber multiplicity alone therefore establishes neither D2 nor its impossibility under injectivity.
    \end{remark}

  \subsubsection*{Bounded flux and the infrared action: an open selection proposal}
    \label{subsec:lemma3}

    \begin{remark}[Status of the saturation proposal]
      \label{lem:saturation-selects}
      The bound $|\partial_t\chi|\leq c_{\mathrm{BI}}$ is a modelling input in the proposed relaxation description.
      Its necessity does not follow from D1--D4, Hypothesis~H, or non-injectivity.
      Nor do those assumptions prove that it selects a Born--Infeld gravitational action or fixes its coefficients.
      No implication between an unbounded flux gradient and collapse of a projection fiber is established here.
    \end{remark}

    In four dimensions, for a local, diffeomorphism-invariant pure-metric action admitting a derivative expansion
    about small curvature, the zero- and two-derivative bulk terms are the cosmological and Einstein--Hilbert terms,
    up to boundary terms.
    Higher-curvature operators enter at higher derivative order with independent coefficients~\cite{Donoghue1994EFT}.
    This ordering requires a regime in which the derivative expansion is controlled; D1--D4 alone do not provide it.
    It requires no flux bound, and does not determine the coefficients or guarantee a nonzero coefficient of $R$.
    Separately, Lovelock's classification restricts a local, symmetric, divergence-free metric tensor depending
    only on the metric and its first two derivatives in four dimensions to a constant linear combination of
    the Einstein tensor and the metric~\cite{Lovelock1971,Lovelock1972}.
    This is a restriction on second-order pure-metric equations, not a selection of a nonlinear completion.

    A Born--Infeld gravitational completion remains a candidate construction.
    To establish a selection result one would have to specify an admissible class of actions and derive which
    member, if any, the microscopic rule and flux bound select.
    A microscopic relation of the schematic form
    \begin{equation}
      G = G(c_{\mathrm{BI}},\hbar_\chi,\ldots)
      \label{eq:G-prediction2}
    \end{equation}
    is a target of that construction, not a quantitative prediction supplied by this section.
    Its normalization and dimensional consistency also remain to be established.

  \subsubsection*{Proposition: scope of the fiber obstruction}
    \label{subsec:necessity-theorem}

    \begin{proposition}[Non-injectivity is necessary for nontrivial fiber-induced structure]
      \label{thm:necessity}
      For an admissible principal projection bundle as in Definition~\ref{def:admissible-projection}:
      \begin{enumerate}[label=(\alph*)]
        \item \label{thm:a}
        a nontrivial vertical re-identification symmetry requires a non-injective projection;
        \item \label{thm:b}
        nontrivial fiber holonomy or nonzero curvature of its principal connection likewise requires
        a non-injective projection.
      \end{enumerate}
      These are necessary conditions, not sufficient conditions for either structure or for gravity.
    \end{proposition}

    \begin{proof}
      Both statements are contrapositives of Lemma~\ref{lem:injectivity-no-gravity}.
      If $\Pi$ is injective, every vertical symmetry is the identity, and the principal structure group,
      its holonomy and its connection curvature are trivial.
    \end{proof}

    A different necessity statement is supplied by ENI~\cite{Beau2026l}:
    a surjective projection whose observable structure is entirely transported from its fine-grained domain,
    with no independent resolving structure, cannot be injective if the resulting structured space is not
    isomorphic to that domain.
    This citation retains ENI's own hypotheses and notion of structural emergence
    (see Appendix~\ref{subsec:non-injectivity-not-ad-hoc}); it is not a theorem excluding injective emergent gravity
    from D1--D4 and Hypothesis~H.
    ENI's domain $\Omega$ is an abstract fine-grained configuration space;
    identifying it with a configuration space of the proposed substrate $\chi$ is a separate modelling step.

  \subsubsection*{Remaining constructive questions}
    \label{subsec:necessity-open}

    \paragraph{Step A: Fiber classification and its limit.}
      Appendix~\ref{subsec:injective-holonomy} proves the injective principal-bundle conclusions used above.
      It does not identify vertical symmetries with diffeomorphisms of the base and does not exclude
      injective projections as candidates for other mechanisms of effective geometry.

    \paragraph{Step B: A gravitational action from the coherence functional.}
      The proposed integrated cost of projective defect is
    \begin{equation}
      S_{\mathrm{coh}}[g,\psi]
      = \int_{\mathcal{M}} \mathcal{L}_{\mathrm{defect}}(\Pi,F,y)\,\mathrm{d}^4y.
      \label{eq:EH-result}
    \end{equation}
      An explicit defect functional and its relation to a specified spectral operator remain to be constructed,
      as recorded in Appendix~\ref{subsec:IR}.
      One would then need to derive its low-curvature expansion and coefficients, test a proposed nonlinear
      completion, and establish a dimensionally consistent matching to $G$.
      Effective-field-theory ordering alone supplies none of these microscopic identifications.
      Neither a Born--Infeld completion nor Eq.~\eqref{eq:G-prediction2} is established here.

      Non-injective fibers may also support winding sectors and coherent collective phases when the required
      topology and dynamics are supplied; the collective proposal below makes additional assumptions.
